% Body of the Computers & Security version (cas.tex): body.tex without the
% learned-model appendix and shortened to the journal's 12,000-word limit.  Every number comes from results/*.json;
% see docs/RESULTS.md for the script that produced each one.

\section{Introduction}\label{sec:intro}

Defenders cannot remediate every published vulnerability, so they rank them.
Most publicly disclosed vulnerabilities receive an identifier in the Common
Vulnerabilities and Exposures (CVE) list.
Two public resources dominate that ranking. The Exploit Prediction Scoring
System (EPSS) assigns every CVE a daily probability of exploitation activity
in the following 30 days~\cite{jacobs2021epss,jacobs2023epssv3}. The
Known Exploited Vulnerabilities catalogue (KEV) of the US Cybersecurity and
Infrastructure Security Agency (CISA) lists CVEs with confirmed exploitation
in the wild~\cite{cisa2022kev}. EPSS looks forward and KEV looks back. Neither
says how likely it is that a given CVE has \emph{already} been exploited
without anyone recording it.

Mell and Spring proposed the Likely Exploited Vulnerabilities metric (LEV) to
fill that gap~\cite{mell2025lev}. LEV treats each 30-day EPSS score as the
probability of exploitation within that window and combines the windows since
the CVE first appeared:
\begin{equation}\label{eq:lev}
\mathrm{LEV}(v, d_0, d_n) \;=\; 1 - \prod_{i=0}^{\lfloor (d_n - d_0)/30 \rfloor}
\Bigl(1 - \mathrm{epss}(v, d_i)\cdot w(d_i, d_n)\Bigr),
\qquad d_i = d_0 + 30i,
\end{equation}
where $\mathrm{epss}(v,d)$ is the EPSS score of CVE $v$ on day $d$, $d_0$ is
the first day on which EPSS scored $v$, $d_n$ is the analysis date, and the
weight $w$ is 1 for every complete window and the covered fraction of the
last, partial one. Among the uses the white paper proposes for LEV are
flagging high-LEV CVEs that are absent from KEV as probably exploited, and
measuring how complete KEV is. It also states that assessing LEV's performance will require
collaboration with parties that hold exploitation data~\cite{mell2025lev}. We
found no published evaluation (Section~\ref{sec:related}).

The obvious evaluation compares LEV with KEV, but such a comparison cannot
separate two explanations. If LEV assigns 0.99 to a CVE that is not in KEV,
either LEV is wrong or KEV has missed an exploited CVE, and the second
possibility is precisely the one the white paper proposes LEV to detect. A
KEV-only evaluation therefore measures LEV's error and KEV's blindness
together, in unknown proportion. The confound is not specific to LEV: every
prioritization method that is tuned or tested against KEV inherits
it~\cite{shimizu2026chaining,sherif2026kri,motamedi2026threatlens}.

This paper asks four questions.
\begin{description}
\item[RQ1] Does a faithful reimplementation reproduce the white paper, and how
well calibrated is LEV against KEV?
\item[RQ2] How do LEV's inputs behave over time, and what does that behaviour
do to LEV?
\item[RQ3] How much of the gap measured against KEV is KEV's incompleteness
rather than LEV's error?
\item[RQ4] Can LEV be corrected, or used to anticipate the catalogues?
\end{description}

\noindent We make three contributions.
\begin{enumerate}
\item \textbf{An evaluation of LEV and its inputs (RQ1, RQ2).} Our
reimplementation reproduces the white paper's averages for CVEs published since
2021 to within 7\% (Section~\ref{sec:res-repro}). Against KEV, LEV over-predicts by a factor of 17.3 while still ranking
CVEs. We trace two input effects behind the gap. EPSS model releases, chiefly v4,
rescaled scores in a single day, and LEV compounds across these boundaries as
if they were evidence; this accounts for an estimated 27--62\% of LEV's mass.
And EPSS scores rise, rather than decay, for CVEs that no exploitation list
records, so LEV keeps growing instead of saturating.
\item \textbf{An estimate of catalogue incompleteness and its consequence
(RQ3).} We audit candidate public exploitation sources for dependence, find that two
of them contain KEV in full and that one of these also contains a sensor feed
we use, remove each containment by source, and apply capture-recapture
estimation to the four resulting lists under explicit rules for when an estimate is
identified, applied to every list configuration. KEV records at most 32\% of the exploited CVEs that
three evidence-citing lists observe. Where the estimate is identified, among
CVEs with LEV of at least 0.9, LEV is about two times over-confident rather
than the 4.5 times that KEV suggests, and below about 2.2 times when LEV cannot
see CISA listings.
\item \textbf{What can and cannot be fixed (RQ4).} A monotone correction
restores calibration a year ahead against the broader ground truth, but drifts
because the same LEV value means different things at different CVE ages. LEV gives
a strong but low-recall early signal, and most listings happen too soon after
a CVE is first scored for any history-based metric to anticipate.
\end{enumerate}

The code, the pinned data whose licences allow it, and per-CVE \emph{truth
bounds} are released: for each of 54,612 CVEs, a lower bound on, and where
identified an estimate of, the probability that it was exploited, so that any
prioritization method can be re-scored against more than KEV.

\section{Background}\label{sec:background}

\subsection{EPSS}

EPSS is a gradient-boosted model that scores every published CVE each day with
the probability of exploitation activity in the following 30
days~\cite{jacobs2021epss,jacobs2023epssv3}. Its ground-truth labels come from
intrusion-detection and honeypot (decoy systems that record attack attempts)
telemetry supplied by Fortinet, AlienVault
OTX, the Shadowserver Foundation and GreyNoise~\cite{jacobs2023epssv3}. The v3
model uses 1,477 input features. Among them are the availability of public
exploit code (Exploit-DB, GitHub, Metasploit), offensive-tool templates such as
Nuclei's, and membership of public exploitation lists, including
KEV~\cite{jacobs2023epssv3}.

Two properties of EPSS matter for what follows. First, KEV is an \emph{input}
to EPSS, so a CVE's listing in KEV can raise its later EPSS scores, and any
metric built from EPSS can partly reflect the catalogue it is then judged
against. The white paper's derivation of LEV treats EPSS as blind to past
exploitation~\cite[Sec.~5.1]{mell2025lev}; that holds for the telemetry EPSS is
trained to predict, but not for KEV membership. Second, the Forum of Incident Response and Security Teams (FIRST),
which publishes EPSS, periodically releases new model versions: v3 on
7 March 2023, v4 on 17 March 2025 and v5 on 15 June
2026~\cite{epss2026releases}. Each release can move every score in the corpus
at once.

\subsection{LEV}

LEV is Equation~\eqref{eq:lev}, applied from the CVE's first EPSS date $d_0$ to
an analysis date $d_n$. Each 30-day window takes the EPSS score on its first
day. The white paper also defines a variant, LEV2, which uses every daily score,
each converted to a one-day probability:
\begin{equation}\label{eq:lev2}
\mathrm{LEV2}(v, d_0, d_n) \;=\; 1 - \prod_{d=d_0}^{d_n}
\Bigl(1 - \tfrac{1}{30}\,\mathrm{epss}(v, d)\Bigr).
\end{equation}
The white paper states both equations as inequalities: it argues that EPSS
under-scores CVEs that have already been exploited, so that the right-hand
side is a \emph{lower bound} on the probability of past
exploitation~\cite[Sec.~5]{mell2025lev}. We evaluate the right-hand side, the
value that an implementation reports.

Because LEV is one minus a product of factors no greater than one, it cannot
decrease as $d_n$ advances. The white paper argues that its growth should
nevertheless stall, because EPSS scores for old CVEs that are not being
exploited decay towards zero~\cite[Secs.~6 and~7.2]{mell2025lev}. Section~\ref{sec:res-growth}
tests that premise.

\subsection{KEV}

CISA maintains KEV under Binding Operational Directive (BOD)
22-01~\cite{cisa2022kev}. An entry requires reliable evidence of active
exploitation, an assigned CVE identifier and clear remediation guidance. Its
\texttt{dateAdded} field records when CISA catalogued the CVE, not when
exploitation began. KEV is selective by design: it serves federal remediation
priorities rather than aiming at complete coverage, and industry analyses
already report that it lists a minority of exploited
CVEs~\cite{roytman2026kevparadox,garrity2026state,beardsley2025kevology}. What
has been missing is an estimate of that incompleteness built from freely available data,
and an account of what it does to evaluations that treat KEV as ground truth.

\section{Data}\label{sec:data}

Table~\ref{tab:data} summarizes every source; we describe each in turn. All
retrieval dates and hashes are recorded in the repository.

\begin{table}[t]
\centering\small
\caption{Data sources. ``Dated $\le d_n$'' counts the entries with a first
exploitation date on or before the analysis date $d_n$ = 2026-08-01, which is
what enters the analysis.}\label{tab:data}
\begin{tabularx}{\linewidth}{@{}p{2.1cm}p{3.6cm}Xp{2.6cm}@{}}
\toprule
source & snapshot & size & role \\
\midrule
EPSS & 2021-04-14 to 2026-08-01 & 67 monthly and 185 daily files; 14,456,555 scores for 367,311 CVEs in the monthly files & LEV input \\
CISA KEV & catalogVersion 2026.08.25 & 1,676 CVEs (1,665 ever scored by EPSS; 1,656 dated $\le d_n$) & list CISA \\
CVE List v5 & \texttt{cve\_2026-08-25\_1700Z} & 192,054 records from 2022 onward & metadata \\
CIRCL sightings & retrieved 2026-09-28 & 68,191 ``exploited'' sightings; 969 CVEs from Shadowserver honeypots (911 dated $\le d_n$), 2,396 from Telegram and similar channels (1,957) & lists SENSOR, SOCIAL \\
VulnCheck KEV & retrieved 2026-09-28 & 5,371 records; 4,815 CVEs after source removal (4,590 dated $\le d_n$); 2,912 under the strict rule (2,824) & list VULNCHK \\
ENISA EUVD & retrieved 2026-09-28 & 66 EU KEV CVEs & excluded: derivative \\
\bottomrule
\end{tabularx}
\end{table}

\paragraph{EPSS.} We use 67 monthly snapshots, the first of each month plus the
dates on which Table~5 of the white paper anchors, from 14 April 2021 to
1 August 2026: 14,456,555 scores for 367,311 CVEs. One monthly file
(2024-12-01) is unavailable from the archive and is replaced by the next day's.
A further 185 daily files serve three checks: they locate model releases to the
day, measure each release's one-day step, and bound the error introduced by
monthly sampling (Section~\ref{sec:threats}).

\paragraph{CISA KEV.} We use catalogVersion 2026.08.25, with 1,676 entries.
Of these, 1,665 have EPSS data; the other 11 were never scored and cannot
receive a LEV value.

\paragraph{Exploitation lists.} CIRCL's Vulnerability-Lookup service publishes
``exploited'' sightings together with the source of each~\cite{circl2026lookup};
we retrieved all 68,191 on 28 September 2026. VulnCheck's KEV lists 5,371
records, each with the URL and date of every exploitation report behind
it~\cite{vulncheck2026kev}; we retrieved it with a free community token on the
same day. We also retrieved the known-exploited dump of ENISA's EU
Vulnerability Database (EUVD)~\cite{enisa2026euvd}.
Section~\ref{sec:m-cr} describes how each source became, or did not become, a
list in the analysis.

\paragraph{Analysis population.} Unless stated otherwise, LEV is computed at
$d_n$ = 1 August 2026. The \emph{calibration cohort} consists of the 105,250
CVEs first scored in the EPSS v3 era (from 7 March 2023) and at least six
months before $d_n$, so that each had time to reach KEV. We exclude CVEs that
EPSS backfilled. When EPSS v2 launched on 4 February 2022 its coverage jumped
from 80,800 to 168,325 CVEs, most of them old, and using the first EPSS date as
a proxy for publication would count them as new. A CVE is therefore kept only
if EPSS first scored it within one calendar year of its CVE-ID year. Of the
cohort, 495 CVEs are in KEV, a base rate of 0.47\%.

\paragraph{Which LEV, which population.} The paper uses two LEV variants, as
specified and version-normalized (Section~\ref{sec:m-version}), and three
populations. Table~\ref{tab:ratios} lists every LEV-to-KEV ratio that appears
later so that they can be compared. The capture-recapture population is younger
than the calibration cohort, so its LEV is lower, and the prospective test set
is younger still. Analyses of RQ3 and RQ4 use version-normalized LEV;
Section~\ref{sec:res-cr} also gives the as-specified values.

\begin{table}[t]
\centering\small
\caption{Every ratio of mean LEV to the KEV rate reported in the paper.}\label{tab:ratios}
\begin{tabular}{@{}lrlr@{}}
\toprule
population & $N$ & LEV variant & mean LEV / KEV rate \\
\midrule
calibration cohort (\S\ref{sec:res-calib}) & 105,250 & as specified & 17.3\x \\
calibration cohort & 105,250 & version-normalized & 12.6\x \\
capture-recapture population (\S\ref{sec:res-cr}) & 54,612 & as specified & 11.2\x \\
capture-recapture population & 54,612 & version-normalized & 7.9\x \\
prospective test set (\S\ref{sec:res-correct}) & 44,316 & version-normalized & 6.3\x \\
\bottomrule
\end{tabular}
\end{table}

\paragraph{Provenance.} Every fixed input is recorded with its SHA-256 digest
and checked when it is loaded or downloaded; a mismatch stops the run. The
exception is the VulnCheck snapshot, which cannot be redistributed and whose
mismatch only raises a warning.

\section{Methods}\label{sec:methods}

\subsection{Reimplementation and reproduction}\label{sec:m-repro}

We implement LEV as specified, including the white paper's inclusive day count
for the final partial window, and LEV2 with its division of each score by 30;
because our snapshots are monthly, each snapshot's score stands in for every
day until the next snapshot. When a window starts on a date for which we hold no snapshot, we take
the next available snapshot, the rule the white paper gives for its own
code~\cite[Sec.~10.3]{mell2025lev}.

To check the implementation, we reproduce Table~5 of the white paper, which
gives, for 16 January 2025, the number and the average LEV of the CVEs
published since each of three dates, and of all CVEs. We compute the same
quantities at $d_n$ = 16 January 2025 from the snapshots dated on or before
it, and compare the CVE counts as well as the averages. The white paper
selects CVEs by publication date; we use the first EPSS date as a proxy, with
the backfill filter of Section~\ref{sec:data}. Because LEV grows with $d_n$,
we also report the reproduction error at every monthly $d_n$ from September
2024 to September 2025, with and without the filter.

\subsection{Calibration}\label{sec:m-calib}

We bin the calibration cohort into ten equal-width LEV buckets and compare each
bucket's mean LEV with the fraction of its CVEs that are in KEV, using Wilson
score intervals~\cite{wilson1927} for the observed rates. We summarize
miscalibration by the \emph{over-prediction ratio}
\begin{equation}\label{eq:ratio}
R \;=\; \frac{\bar{L}}{\hat{p}},
\end{equation}
the mean LEV $\bar{L}$ of a group of CVEs divided by the fraction $\hat{p}$ of
them that are exploited under a stated notion of truth, with a 95\% interval
from Poisson error on the count of exploited CVEs. When $\hat{p}$ is KEV
membership we say LEV \emph{over-predicts} KEV; when it is an estimate of true
exploitation we say LEV is \emph{over-confident}. We also report the Brier
score~\cite{brier1950} and the expected calibration error over the same ten
bins~\cite{naeini2015ece},
$\mathrm{ECE} = \sum_b (n_b/N)\,|\bar{L}_b - \hat{p}_b|$. We test for a
monotone trend in the observed rate across buckets with the Cochran--Armitage
test~\cite{armitage1955trend}. Because the comparison could depend on how long a CVE has had to reach
KEV, we repeat it by CVE age at $d_n$ (6--12, 12--24 and more than 24 months)
and for maturity rules of 3, 6 and 12 months.

\subsection{Version normalization}\label{sec:m-version}

EPSS releases can move every score at once, and LEV multiplies scores across
them. We estimate how much of LEV this contributes in two ways.
\emph{Percentile mapping} maps every score produced after the v4 release back
onto the v3 scale: each score keeps its percentile within its own day's file,
and we substitute the score at that percentile in the last v3 file
(16 March 2025). \emph{Step-only} mapping fits the map once at each release,
from the last file before to the first file after (v4 on 17 March 2025, v5 on
15 June 2026), and then holds it fixed. This removes the one-day rescaling and
leaves organic drift within a version untouched.

Neither method can separate a corpus-wide shift from genuine per-CVE
improvement in the newer model, so the mass each removes is an estimate of the
artifact, not a corrected value. We report the two as a range, together with an
upper bound obtained by treating all of the calibration cohort's re-ranking at
the v4 and v5 releases as artifact. We call LEV computed from percentile-mapped scores
\emph{version-normalized} LEV, and use it for RQ3 and RQ4.

\subsection{Analysis-date sweep}\label{sec:m-sweep}

To see how LEV evolves, we follow two fixed cohorts: the 13,685 CVEs first
scored between March and September 2023, and the 7,558 first scored between
May and September 2021. For each we compute LEV at every month from October
2023 to August 2026. We compare linear, logarithmic, two- and three-parameter
saturating, quadratic, exponential and power-law growth. Monthly residuals are
strongly autocorrelated (lag-1 correlation 0.6--0.9), so we judge the comparison
with heteroskedasticity- and autocorrelation-consistent (HAC) standard
errors~\cite{newey1987hac} and a moving-block
bootstrap~\cite{kunsch1989bootstrap} rather than with an information criterion
alone.
The white paper's premise is tested directly: we track the mean EPSS of cohort
CVEs that no exploitation list ever records.

\subsection{Capture-recapture}\label{sec:m-cr}

Capture-recapture, also called multiple-systems estimation, estimates the size
of a population from the overlap between incomplete lists that each record part
of it~\cite{iwgdmf1995,bird2018mse}. If two independent lists each record half
of a population, about a quarter of the population should appear on both, and
the observed overlap therefore reveals how many members both lists missed.

\paragraph{Lists.} An \emph{exploitation list} is a set of CVEs that a source
records as exploited in the wild, each with the date it was first recorded. We
use four. CISA is CISA KEV. SENSOR consists of CIRCL sightings from
Shadowserver's honeypot feed, which records attempted exploitation directly.
VULNCHK is VulnCheck's KEV, rebuilt from its per-report evidence as described
below. SOCIAL consists of CIRCL sightings from Telegram and similar channels.

\paragraph{Dependence audit.} Capture-recapture assumes that the lists are
separate observations. A list that contains another makes their overlap equal
to the smaller list by construction, which would imply that nothing is missing.
We therefore audited every candidate before use, and found three containments.
CIRCL's aggregated KEV catalogue contains all 1,676 CISA KEV CVEs, so we exclude
it. VulnCheck's KEV includes CISA KEV by its own
documentation~\cite{vulncheck2026kev}, so we rebuild it from its reports and
drop every report hosted on \texttt{cisa.gov}. The third was less obvious.
VulnCheck's largest single source is Shadowserver's dashboard, with 401,967
reports, and every one of the 969 SENSOR CVEs appears in VulnCheck's
Shadowserver subset. We drop those reports too. What remains of VulnCheck's KEV
is its own honeypot canaries, SANS DShield, CrowdSec, vendor advisories and
threat reports. For 78\% of the CISA-listed CVEs it contains, the earliest
remaining report predates CISA's listing (median lead 28 days over all of
them), so the remainder is not an echo of CISA.

EU KEV lists 66 CVEs, 55 of which EUVD itself also attributes to CISA KEV. Of
the 41 that appear in our pinned CISA snapshot, 36 (88\%) carry CISA's own
\texttt{dateAdded}, so we treat EU KEV as derivative and exclude it. EUVD's
exploited-flag search returns the same consolidated set. Exploit-DB records the
availability of exploit code, not exploitation, and is also excluded.

\paragraph{Primary lists.} Telegram channels carry mass lists of CVEs that
actors claim to use, and social-media signals are known to be
noisy~\cite{sabottke2015twitter}. We therefore designate CISA, SENSOR and
VULNCHK, the three lists whose entries each cite evidence, as the
\emph{primary} lists, and report every other configuration alongside.

\paragraph{Population.} SENSOR begins on 24 October 2024. We restrict the
analysis to the 54,612 CVEs whose first monthly EPSS snapshot is dated between
24 October 2024 and 30 January 2026 (in practice, the snapshots of 1 November
2024 to 1 January 2026), so that every list could observe nearly every CVE from
its first score, and apply the six-month maturity rule. Because snapshots are
monthly, 2,365 of these CVEs (4\%) were in fact first scored shortly before
SENSOR began: they already appear in EPSS's daily file of 22 October 2024. In all, 734 of the 54,612 CVEs appear on at least
one primary list. Capture probabilities plainly vary between CVEs; a CVE that
attracts wide exploitation is more likely to appear on every list. We absorb
part of this heterogeneity by stratifying on version-normalized LEV into
$[0, 0.1)$, $[0.1, 0.9)$ and $[0.9, 1]$, and fitting each stratum separately.

\paragraph{Models.} Two-list estimates computed from pairs of the four lists
disagree by a factor of 2.4 (530 to 1,256 exploited CVEs), so the lists are not
independent and dependence must be modelled. For $t$ lists, let $n_h$ be the
number of CVEs with capture history $h \in \{0,1\}^t$, $h \neq 0$. We fit
Poisson log-linear models~\cite{bishop1975,iwgdmf1995}
\begin{equation}\label{eq:loglin}
\log \mathrm{E}[n_h] \;=\; \lambda + \sum_{j} \lambda_j h_j
+ \sum_{(j,k) \in I} \lambda_{jk} h_j h_k
\end{equation}
to the $2^t - 1$ observable patterns under every hierarchical set
$I$ of two-way interactions. The number of CVEs on no list is
estimated as $\hat{n}_0 = \exp(\hat\lambda)$, which requires the standard
identifying assumption that there is no $t$-way
interaction~\cite{iwgdmf1995,bishop1975}. We select among models by AIC and
obtain a 95\% interval by parametric bootstrap from the fitted cell means
(1,000 replicates). A model whose interaction term is fitted to an empty cell
reaches its maximum likelihood only at the boundary of the parameter space; we
score such a model by its extended maximum-likelihood estimate and count its
residual degrees of freedom accordingly~\cite{fienberg2012mle}, and exclude it
as saturated if none remain.

Chao's estimator~\cite{chao1987} gives a lower bound that remains valid when
capture probabilities differ between CVEs, though not when lists are negatively
dependent (Section~\ref{sec:res-cr} checks the one pair where our construction
could make them so). In its form for $t$ lists~\cite{chao1989sparse} it is
\begin{equation}\label{eq:chao}
\hat{N}_{\mathrm{Chao}} \;=\; S + \frac{t-1}{t}\cdot\frac{f_1^2}{2 f_2},
\end{equation}
where $S$ is the number of CVEs on at least one list and $f_1$ and $f_2$ are the
numbers on exactly one and exactly two lists. We report it with a log-normal
interval.

\paragraph{When not to report an estimate.} Log-linear estimates of the unseen
cell can be driven entirely by the choice of interaction terms, especially when
lists do not overlap~\cite{chan2021sparse,binette2022mse}. We call a
population \emph{identified} only if every model within 2 AIC units of the best,
the set that Burnham and Anderson describe as having substantial
support~\cite{burnham2002}, (i)~stays within the population, (ii)~does not fall
below the lower end of Chao's interval, and (iii)~for a pooled population, does
not fall below the sum, over the three LEV strata, of the lower ends of their
Chao intervals (the observed count where $f_2 = 0$); and if, in addition, those
models agree within a factor of two, a criterion of our own. Otherwise we report no point estimate. The
rules were not registered in advance. Rules (i) and (ii) and the factor-of-two
condition were written together with the first estimates; rule (iii) was added
later, when pooled estimates were seen to hide the heterogeneity between
strata. Because each rule is an additional condition, adding one can only
withhold an estimate, never create one. Table~\ref{tab:config} applies the
same rules to every list configuration.

\paragraph{The truth ladder.} Rather than a single estimate, we report LEV's
over-confidence against successively stronger notions of truth, each adding one
assumption to the one before: KEV alone; KEV or SENSOR; the union of the
primary lists; Chao's bound; and the log-linear estimate. Every rung is a lower
bound on, or an estimate of, the number of exploited CVEs, so every
over-confidence figure is an upper bound on, or an estimate of, LEV's error
under that rung's assumption. A reader can stop at whichever assumption they are willing to
accept.

\subsection{Correction curve}\label{sec:m-correct}

We fit monotone maps from LEV to observed exploitation by isotonic
regression~\cite{zadrozny2002}, by Platt scaling~\cite{platt1999}, and by a
Platt variant that also takes log CVE age as an input. There are two targets:
KEV on the calibration cohort, and the union of the primary lists on the
capture-recapture population. All scores are computed on held-out CVEs under
three schemes. Five-fold cross-validation is the in-distribution case. A split
by first-score date trains on older CVEs and tests on newer ones. The deciding
scheme is prospective: the map is fitted at $d_n$ = 1 August 2025 using the
outcomes known then, and applied at 1 August 2026 to the 44,316 CVEs first
scored after the training cohort, none of which it saw. Prospective ratios carry
bootstrap intervals over test CVEs.

\subsection{Circularity and early signal}\label{sec:m-leak}

Because KEV is an EPSS input, we measure how much listing moves scores, using
two designs that cannot select on the outcome.

\paragraph{The jump at listing.} For each KEV CVE we compare its EPSS score in
the last snapshot before \texttt{dateAdded} with the first snapshot on or after
it. We net this change against non-KEV CVEs in the same starting-score decile
over the same pair of snapshots, and exclude pairs that straddle a model
release, leaving 1,168 KEV CVEs.

\paragraph{Pre-listing LEV.} We compute each KEV member's LEV using only
snapshots dated strictly before its listing, and compare it with matched
controls: non-KEV CVEs first scored on the same date and evaluated on the same
date. As a forward check, we ask how well LEV of 0.5 or more at a fixed CVE age
(one, three, six or twelve months) predicts later listing.

\section{Results}\label{sec:results}

\subsection{RQ1: reproducing the white paper}\label{sec:res-repro}

Table~\ref{tab:repro} compares our CVE counts and average LEV with Table~5 of
the white paper, at the date the white paper gives for it. The three rows that
are defined by a publication date agree to within 7\%. Our cohorts hold 7--8\%
fewer CVEs, as expected: the first monthly EPSS snapshot stands in for the
publication date, so CVEs published in the two weeks before $d_n$ are not yet
counted, and the backfill filter of Section~\ref{sec:data} removes some
genuine late publications. We added that filter after the two older rows first
missed by a factor of more than two, but it is justified independently: the
first EPSS date is a proxy for publication only for CVEs that EPSS did not
backfill.

\begin{table}[t]
\centering\small
\caption{Table 5 of the white paper~\cite{mell2025lev} against our
reimplementation at the same date, $d_n$ = 2025-01-16: number of CVEs and
their average LEV.}\label{tab:repro}
\begin{tabular}{@{}lrrrrr@{}}
\toprule
 & \multicolumn{2}{c}{CVEs} & \multicolumn{3}{c}{average LEV} \\
\cmidrule(lr){2-3}\cmidrule(l){4-6}
CVEs published since & white paper & ours & white paper & ours & difference \\
\midrule
2023-03-07 & 66,125 & 60,934 & 0.015 & 0.0160 & $+7\%$ \\
2022-02-04 & 94,440 & 87,866 & 0.047 & 0.0495 & $+5\%$ \\
2021-04-14 & 112,754 & 103,718 & 0.087 & 0.0873 & $+0.3\%$ \\
all CVEs & 273,979 & 272,702 & 0.146 & 0.2593 & $+78\%$ \\
\bottomrule
\end{tabular}
\end{table}

The match does not depend on the exact date. With the filter, the three rows
are within 12\% at each monthly $d_n$ from September 2024 to March 2025. From
April 2025, after the v4 release, the error grows, and by September the worst
row is off by 200\%. Without the filter, the two older rows are off by more
than 140\% at every date.

The fourth row, the average over every CVE that EPSS had scored, is not
reproduced. Our population matches the white paper's to within 0.5\%, but our
average is 78\% higher. The difference lies in the roughly 164,000 CVEs with
identifiers from 2020 or earlier, which EPSS first scored at its launch or
backfilled later. The white paper notes that such CVEs ``are missing many EPSS
scores''~\cite[Sec.~7.2]{mell2025lev}, which lowers LEV, and we cannot tell
which scores its computation lacked; dropping all scores from before EPSS v2,
or from before v3, does not recover its value either. Restricting the row to
CVEs that EPSS did not backfill gives 0.135 ($-7\%$), but over fewer than half
as many CVEs (123,163), so we do not count it as a match. No LEV value
reported below is computed for these older CVEs: every cohort consists of CVEs
that EPSS scored when they were published, the kind in the first three rows. The
reproduction validates the implementation; it does not validate the metric.

\finding{A faithful reimplementation reproduces the white paper's average LEV
for CVEs published since 2021 to within 7\% at the white paper's own analysis
date. Its average over all CVEs, which is dominated by older CVEs with
incomplete EPSS histories, is not reproduced.}

\subsection{RQ1: against KEV, LEV over-predicts 17-fold but ranks well}\label{sec:res-calib}

Across the 105,250-CVE calibration cohort, mean LEV is 0.081 against a KEV
rate of 0.0047: LEV over-predicts KEV by a factor of 17.3 (95\% interval
15.8--18.9, from Poisson error on the 495 KEV CVEs). The Brier score is 0.0373
and the ECE 0.0765.

LEV nevertheless ranks. The observed KEV rate climbs from 0.03\% in the lowest
LEV bucket to 12.9\% in the highest, where LEV averages 0.987
(Figure~\ref{fig:reliability}). The trend is not carried by the extremes alone:
within the eight middle buckets it remains strong (Cochran--Armitage $z$ = 9.2).

\begin{figure}[t]
\centering
\fig{reliability.pdf}
\caption{Reliability of LEV against CISA KEV on the calibration cohort, on log
scales. Error bars are 95\% Wilson intervals; dotted lines mark 10\x{} and
100\x{} over-prediction. The lower panel gives the number of CVEs in each
bucket.}\label{fig:reliability}
\end{figure}

Over-prediction grows with CVE age (Table~\ref{tab:age}), the first sign of an
effect that recurs in Sections~\ref{sec:res-growth} and~\ref{sec:res-correct}.
The six-month maturity rule does not drive the result: with three or twelve
months the ratio is 16.0 or 19.1.

\begin{table}[t]
\centering\small
\caption{Calibration against KEV by CVE age at $d_n$. Intervals are
95\%.}\label{tab:age}
\begin{tabular}{@{}lrrrr@{}}
\toprule
CVE age at $d_n$ & $N$ & KEV CVEs & mean LEV / KEV rate & KEV rate at LEV $\ge$ 0.9 \\
\midrule
6--12 months & 22,043 & 82 & 6.6\x{} [5.3, 8.3] & 27.1\% \\
12--24 months & 40,875 & 200 & 15.1\x{} [13.1, 17.4] & 16.3\% \\
24+ months & 42,332 & 213 & 23.4\x{} [20.5, 26.9] & 9.8\% \\
\bottomrule
\end{tabular}
\end{table}

\finding{Against KEV, LEV over-predicts by a factor of 17.3 but ranks CVEs
across its whole range. Over-prediction is three to four times larger for CVEs
older than two years than for CVEs less than a year old.}

\subsection{RQ2: EPSS releases put an artifact into LEV}\label{sec:res-version}

EPSS v4 went live on 17 March 2025~\cite{epss2026releases}. Daily files locate
it exactly. On an unchanged set of 13,552 CVEs from the 2023 sweep cohort, mean
EPSS goes from 0.00862 on 16 March 2025 to 0.02640 the next day, a 3.06-fold
step (Figure~\ref{fig:drift}). The whole corpus moves by a factor of only 1.21,
and the cohort's Spearman rank correlation across the step is 0.61: for young
CVEs, v4 re-ranked as much as it rescaled. Of the rise in cohort mean EPSS over
2023--2026, 64\% (2023 cohort) and 70\% (2021 cohort) happens at release
boundaries rather than between them. EPSS v5 (15 June 2026) moved the same
cohort's mean the other way, to 0.79 of the previous day's value, while its
median more than doubled (from 0.0021 to 0.0054).

\begin{figure}[t]
\centering
\fig{epss_drift.pdf}
\caption{Cohort mean EPSS, monthly, with daily detail around the v4 and v5
releases.}\label{fig:drift}
\end{figure}

LEV multiplies scores across these boundaries as if they belonged to one
measurement series. Mapping post-v4 scores back onto the v3 scale lowers mean
LEV on the calibration cohort by 26.9\% under percentile mapping, which also
removes drift within each version; removing only the one-day steps (step-only
mapping) lowers it by 29.0\%, almost all of it at v4 (28.6\% at v4 alone,
3.2\% at v5 alone). Over-prediction against KEV falls from 17.3 to between 12.3
and 12.6; version-normalized LEV has a Brier score of 0.0289 and an ECE of
0.0547 (0.0281 and 0.0530 under step-only mapping). Treating \emph{all} of the cohort's re-ranking at the v4 and v5
releases as artifact
removes 62\% of LEV's mass, which we take as an upper bound.

\finding{Model releases explain an estimated 27--62\% of LEV's
mass. Even at the upper bound, LEV still over-predicts KEV by a factor of
6.6.}

\subsection{RQ2: LEV grows without saturating, because EPSS does not decay}\label{sec:res-growth}

For the 2023 cohort, mean LEV rises from 0.008 in October 2023 to 0.135 in
August 2026; for the 2021 cohort it rises from 0.31 to 0.41 over the same
period (Figure~\ref{fig:sweep}). Monotone growth is a property of
Equation~\eqref{eq:lev}, not a finding. The empirical questions are whether the
growth slows, and whether the premise behind an expected slowdown holds.

It does not slow. For the 2023 cohort, a three-parameter saturating model
collapses to its linear limit. Within the stretch with no model release, the
curvature term is significant under autocorrelation-robust errors (HAC
$z$ = 3.33), and a moving-block bootstrap rejects linear ($p$ = 0.034) and
saturating ($p$ = 0.018) growth in favour of accelerating growth. Within that
stretch a quadratic fits clearly better than a power law; over the full series,
which includes the v4 step, the two fit almost equally well.

\begin{figure}[t]
\centering
\fig{dn_sweep.pdf}
\caption{Mean LEV of two fixed cohorts as the analysis date advances. Dotted
lines mark EPSS releases. Further panels are in Appendix~\ref{app:robust}.}\label{fig:sweep}
\end{figure}

The premise fails directly. For the 13,390 CVEs in the 2023 cohort that no
exploitation list records, mean EPSS rises from 0.0022 to 0.0046 within the
single v3 model (monthly trend $p < 10^{-7}$ with ordinary and with
autocorrelation-robust standard errors), and reaches 0.0170
by August 2026 across releases. Scores for old, unlisted CVEs rise; they do not
decay. Without decay, LEV keeps compounding, and over time more of the corpus
approaches one.

\finding{LEV does not saturate. The decay of EPSS for unexploited CVEs that
the white paper relies on is not observed: EPSS for never-listed CVEs roughly
doubles within one model version.}

\subsection{RQ3: how much of the gap is the catalogue's}\label{sec:res-cr}

Table~\ref{tab:ladder} gives the truth ladder on the 54,612-CVE
capture-recapture population with the three primary lists, pooled over LEV
strata.

\begin{table}[t]
\centering\small
\caption{LEV over-confidence against successively stronger notions of truth,
capture-recapture population, primary lists, all LEV strata pooled. Intervals
are 95\%: Poisson on the counts, log-normal for Chao's bound.}\label{tab:ladder}
\setlength{\tabcolsep}{3.5pt}
\begin{tabular}{@{}llrrr@{}}
\toprule
truth & assumption added & exploited & version-normalized LEV & as-specified LEV \\
\midrule
CISA KEV & KEV is complete & 236 & 7.93\x{} [6.98, 9.05] & 11.16\x \\
CISA or SENSOR & honeypot sightings are real & 336 & 5.57\x{} [5.01, 6.22] & 7.84\x \\
union of primary lists & every listed report is real & 734 & 2.55\x{} [2.37, 2.75] & 3.59\x \\
Chao lower bound & closed population & 974 & 1.92\x{} [1.79, 2.04] & 2.70\x \\
log-linear & no three-way interaction & \multicolumn{3}{c}{not identified} \\
\bottomrule
\end{tabular}
\end{table}

By counting alone, LEV's over-confidence falls from 7.9 to 2.6. Of the CVEs
that any primary list records, 68\% [64, 71] are absent from KEV, or 60\%
[55, 64] under the strict VulnCheck rule defined below. KEV
recorded at most 32\% of the exploited CVEs that the primary lists observed,
and at most 24\% of Chao's lower bound. The pooled log-linear estimate (881)
falls below Chao's bound (919 at the lower end of its interval) and is
therefore not reported. Table~\ref{tab:strata} stratifies by LEV.

\begin{table}[t]
\centering\small
\caption{Over-confidence by LEV stratum, primary lists, version-normalized
LEV. Intervals are 95\%.}\label{tab:strata}
\begin{tabular}{@{}lrrrrrl@{}}
\toprule
LEV & $N$ & mean LEV & vs KEV & vs union & vs Chao & log-linear \\
\midrule
$[0.0, 0.1)$ & 52,012 & 0.012 & 21.2\x & 2.36\x & 1.05\x{} [0.83, 1.29] & not identified \\
$[0.1, 0.9)$ & 1,964 & 0.322 & 9.16\x & 3.21\x & 2.54\x{} [2.24, 2.76] & not identified \\
$[0.9, 1.0]$ & 636 & 0.985 & 4.54\x & 2.26\x & 2.05\x{} [1.93, 2.14] & \textbf{2.04\x{} [1.81, 2.32]} \\
\bottomrule
\end{tabular}
\end{table}

\paragraph{The top stratum is identified.} There, LEV averages 0.985, and we
estimate that 48\% [42, 54] of those CVEs were exploited: LEV is about two times
over-confident, against the 4.5 times that KEV alone suggests. As-specified LEV
gives the same figure in this stratum (2.04).

\paragraph{The middle stratum is not identified.} No CVE in it is on both CISA
and SENSOR. The two best models fit that empty cell at the boundary and have no
residual degrees of freedom, and the model that remains falls below Chao's
bound. Chao's bound alone caps over-confidence there at 2.5 (3.2 by counting).

\paragraph{The lowest stratum is not identified either,} and it holds 95\% of
the CVEs. Every supported model falls below Chao's bound. Chao's bound alone
brings over-confidence to 1.05 [0.83, 1.29]: in this stratum, LEV may not
over-predict at all.

\begin{figure}[t]
\centering
\fig{capture_recapture.pdf}
\caption{Exploitation rate under each notion of truth against mean LEV, by LEV
stratum, primary lists. The diagonal is perfect calibration. Error bars are
95\% intervals.}\label{fig:cr}
\end{figure}

\paragraph{Does this depend on the choice of lists?} Table~\ref{tab:config}
repeats the log-linear estimate under every list configuration. Wherever the top
stratum is identified, its estimate lies between 1.92 and 2.12. That includes
the strict configuration, which keeps a VulnCheck CVE only if a sensor saw it
or at least two independent sources reported it (2,912 CVEs). This check
matters, because VulnCheck supplies 398 of the 430 CVEs that appear on a
single primary list, and 38\% of those 398 rest on reports from one source with
no sensor evidence.

\begin{table}[t]
\centering\small
\caption{Log-linear over-confidence by list configuration. ``Not claimed'':
the pooled population passes the identification rules, but it pools strata
that the stratified analysis rejects.}\label{tab:config}
\begin{tabular}{@{}llll@{}}
\toprule
lists & LEV $[0.9, 1.0]$ & LEV $[0.1, 0.9)$ & all LEV \\
\midrule
CISA, SENSOR, VULNCHK (primary) & 2.04\x{} [1.81, 2.32] & not identified & not identified \\
CISA, SENSOR, strict VULNCHK & 2.12\x{} [1.86, 2.45] & not identified & not claimed (2.64\x) \\
CISA, SENSOR, SOCIAL & 2.07\x{} [1.75, 2.44] & not identified & not claimed (0.99\x) \\
CISA, SOCIAL, VULNCHK (no SENSOR) & 1.92\x{} [1.64, 2.25] & 1.22\x{} [0.67, 1.72] & not identified \\
all four & not identified & not identified & not claimed (0.35\x) \\
\bottomrule
\end{tabular}
\end{table}

Three further checks bear on the top-stratum result. First, Shadowserver is
one of EPSS's label sources, so EPSS is trained to predict roughly what SENSOR
records. Counting with CISA and VULNCHK alone gives 2.65 pooled and 2.38 in the
top stratum. Second, removing Shadowserver-sourced reports from VulnCheck may
create negative dependence between SENSOR and VULNCHK by construction. Forcing
that interaction into every model leaves no population identified, and gives
2.24 in the top stratum.

Third, KEV is an EPSS input, so a CISA listing can lift a CVE into the top
stratum (Section~\ref{sec:res-leak}). We therefore re-formed the strata on LEV
that cannot see a CISA listing, in two ways. Freezing each CISA-listed CVE's LEV
on the day before its listing (133 of the 236 had no earlier EPSS score and get
LEV 0) leaves the top stratum at 2.16 against Chao's bound and 3.43 against the
union, compared with 2.05 and 2.26 above. This design favours unlisted CVEs,
which keep accumulating windows, so we also measured every CVE's LEV at the
same age, 90 or 180 days after its first score, and dropped the CVEs that CISA
had listed by then. The top stratum, now 50 and 172 CVEs, is 1.04 and 1.66
times over-confident against Chao's bound, and 1.37 and 2.43 times against the
union. The log-linear estimate is not identified in any of these variants, but
since Chao's bound is a lower bound on the exploited count, over-confidence in
the top stratum stays below about 2.2 in all of them. Against KEV alone the
picture changes: once LEV cannot see the listing, over-prediction of KEV in the
top stratum rises from 4.5 to between 12 and 99. A KEV-only evaluation of
high-LEV CVEs is thus flattered by the very feedback it ignores. Freezing LEV
at the first listing on \emph{any} primary list is harsher (6.8 against Chao's
bound and 8.3 by the log-linear estimate, which is identified here, with 26
listed CVEs among 385), but it also discards SENSOR and VULNCHK
sightings, which do not feed EPSS's per-CVE inputs, and so tests anticipation
rather than circularity; Section~\ref{sec:res-leak} shows that LEV rarely
anticipates listings.

\finding{KEV records at most a third of the exploited CVEs that three
evidence-citing lists observe. Where the estimate is identified, LEV$\,\ge\,$0.9
is about two times over-confident, not 4.5 times, and stays below about 2.2
times when LEV cannot see CISA listings. Below that, the data bound LEV's error
but do not identify it.}

\subsection{RQ4: a correction curve, and why it drifts}\label{sec:res-correct}

Table~\ref{tab:prosp} gives the prospective validation against KEV. Under
cross-validation, any of the three maps restores calibration (ratio of
predicted to observed 1.01). Prospectively against KEV, each map cuts the
calibration error (ECE) of version-normalized LEV by a factor of 12 to 17, but all of them under-predict by
28--39\% a year later.

\begin{table}[t]
\centering\small
\caption{Prospective validation against KEV: fitted at 2025-08-01 and scored
at 2026-08-01 on 44,316 CVEs the fit never saw (192 in KEV).}\label{tab:prosp}
\begin{tabular}{@{}lrrl@{}}
\toprule
method & Brier & ECE & predicted / observed \\
\midrule
LEV (version-normalized) & 0.0113 & 0.0230 & 6.30\x{} [5.56, 7.27] \\
LEV (as specified) & 0.0136 & 0.0339 & 8.81\x{} [7.77, 10.18] \\
isotonic & 0.0036 & 0.0020 & 0.61\x{} [0.53, 0.70] \\
Platt & 0.0037 & 0.0015 & 0.66\x{} [0.58, 0.76] \\
Platt with CVE age & 0.0036 & 0.0014 & 0.72\x{} [0.63, 0.83] \\
\bottomrule
\end{tabular}
\end{table}

Against the union of the primary lists, the maps are fitted on 111 training
positives and scored on 597 test positives. The isotonic map gives a ratio of
1.00 [0.93, 1.08], Platt scaling 0.80 [0.74, 0.87] and age-aware Platt scaling
0.88 [0.81, 0.95]. Much of the drift against KEV therefore reflects KEV's
narrow and lagging coverage rather than a failure of the correction itself.

Transfer across birth cohorts is worse. A map fitted on CVEs first scored
before 1 June 2025 under-predicts younger CVEs at the same analysis date
(isotonic 0.30, Platt 0.52 against the union), and the age-aware Platt model
overshoots (2.32). The drift has a cause we can name. An old CVE reaches a given
LEV through many windows of modest EPSS, a new one through a few high scores.
The second is stronger evidence, so the same LEV value carries a different
probability at different ages. This is Table~\ref{tab:age} seen from another
side. Adding log age to Platt scaling narrows the prospective shortfall
against KEV only from 34\% to 28\% (Table~\ref{tab:prosp}).

Table~\ref{tab:curve} gives the isotonic correction fitted on all data, which
is the form a practitioner would use. Both targets are catalogue appearances,
so both are lower bounds on true exploitation. The intervals are several points
wide: the curve conveys the scale of LEV's error, not a precise mapping for a
single score.

\begin{table}[t]
\centering\small
\caption{The isotonic correction, fitted on all data, with 95\% bootstrap
intervals (400 resamples of CVEs).}\label{tab:curve}
\begin{tabular}{@{}lll@{}}
\toprule
LEV & corrected to KEV & corrected to union of primary lists \\
\midrule
0.10 & 0.8\% [0.5, 1.0] & 6.1\% [4.6, 7.5] \\
0.50 & 2.7\% [1.9, 4.1] & 15.5\% [11.5, 19.0] \\
0.90 & 6.7\% [4.4, 8.4] & 20.2\% [16.5, 29.5] \\
0.99 & 8.2\% [6.4, 10.1] & 32.5\% [26.9, 37.4] \\
1.00 & 27.8\% [24.3, 30.9] & 76.3\% [68.0, 85.9] \\
\bottomrule
\end{tabular}
\end{table}

\begin{figure}[t]
\centering
\fig{recalibration.pdf}
\caption{Left: the isotonic correction curves. Centre and right: held-out ratio
of predicted to observed exploitation under each validation scheme, against KEV
and against the union of primary lists (log scale; 1 is perfect
calibration).}\label{fig:recal}
\end{figure}

\finding{A monotone correction makes LEV's values usable a year ahead against
the broader ground truth, but must be refitted as CVEs age.}

\subsection{RQ4: LEV's early signal is strong but rare, and its top bucket is partly circular}\label{sec:res-leak}

\paragraph{The listing feedback.} At the listing date, the median KEV CVE's
EPSS does not move relative to matched non-KEV controls (median ratio 1.00 over
1,168 CVEs). A quarter of listed CVEs, however, rise at least 2.85 times more
than their controls. The KEV input feature matters for a sizeable minority of
listings rather than for all of them.

\paragraph{Before listing.} Of the cohort's 495 KEV members, 56\% were listed
before EPSS's first monthly snapshot of them, so our snapshots hold no
pre-listing history for them, and even daily scores would give a history-based
metric at most about a month to flag them first. The other 220 have a pre-listing history (median
2.8 months). Using only snapshots dated strictly before listing, 15\% of them
had LEV of 0.5 or more, against 0.58\% of matched controls: a high LEV was 25
times more common among CVEs that KEV would later list. The typical member
already stood out, with a median pre-listing LEV of 0.005, at the 92nd
percentile of its controls. After listing, the members' median LEV is 1.0,
whereas the same controls move only from 0.0013 to 0.017 over the same period,
so elapsed time explains little of that jump.

\paragraph{As a forecaster.} At a fixed CVE age of three months, a LEV of 0.5
or more identifies CVEs that are listed later with 6.0\% precision, 59 times
the base rate, but with only 10\% recall. At six months the precision is 2.4\%
and the recall 16\%. At one month it flags only 31 CVEs (precision 9.7\%,
recall 2\%); at twelve months the precision falls to 0.5\%, 13 times the base
rate, with recall 16\%.

\finding{LEV carries a real early signal, but for few CVEs. Most of the
concentration of KEV members in the top LEV bucket arrives after listing and is
partly circular.}

\section{Discussion}\label{sec:discussion}

\paragraph{What LEV is good for.} LEV ranks. Its extremes separate, with the
top bucket's KEV rate nearly 500 times the bottom bucket's, and the trend holds
across the middle of its range. Its values should not, however, be read as
probabilities, still less as the lower bounds on them that the white paper
derives. A LEV of 0.9 corresponds to about a 7\% chance of KEV listing,
and about a 20\% chance of appearing on any primary list
(Table~\ref{tab:curve}).

\paragraph{What LEV is not yet good for.} The white paper proposes LEV for
finding exploited CVEs that KEV has missed. Our results cut both ways. LEV's
apparent over-confidence is substantially KEV's incompleteness, which is the
premise of the proposal. Before listing, a high LEV is 25 times more common
among CVEs that KEV later lists. But it flags only 15\% of those with any
pre-listing history (32 of 495 members, or 6\%), and most listings happen within a month of a CVE's first EPSS score, before any history
accumulates. LEV can support the proposal at the margin; it cannot carry it.

\paragraph{What would make LEV more reliable.} Five changes follow from the
results.
(1)~Compute LEV within one EPSS model version, or normalize across releases,
since a release otherwise enters LEV as evidence.
(2)~State the analysis date with every published LEV value, since values
computed at different dates are not comparable.
(3)~Stop accumulating windows, or decay them, for CVEs with no recent evidence,
or accept that LEV drifts towards one.
(4)~Read LEV through a correction that is refitted as CVEs age, not at face
value.
(5)~Validate against more than KEV, which is both an EPSS input and incomplete.

\paragraph{KEV coverage, and evaluations built on it.} Empirical Security
reports, from its own exploitation monitoring, that KEV covers at most about 8.7\% of CVEs
with observed exploitation activity~\cite{roytman2026kevparadox}. Our figures
come from freely available data on a younger population: KEV recorded at most 32\% of what
the primary lists observed, and at most 24\% of Chao's bound. The direction
agrees, and the levels differ as one would expect, since three lists see less
than one vendor's exploitation monitoring. Any method tuned or tested on
KEV~\cite{shimizu2026chaining,sherif2026kri,motamedi2026threatlens} measures,
in part, CISA's visibility.

\paragraph{Two definitions of ``exploited''.} Part of the KEV gap is
definitional. EPSS is trained on intrusion-detection and honeypot detections,
so LEV estimates the probability of observed exploitation \emph{activity}. KEV
lists confirmed, actionable exploitation. Against sources closer to EPSS's own
definition, LEV in its top stratum looks about twice too confident rather than
four or five times.

\paragraph{Implications for policy and practice.} KEV's selectivity is
deliberate. Under BOD 22-01 every entry imposes a remediation deadline on
federal agencies, so each must be confirmed and actionable~\cite{cisa2022kev}.
The problem is not KEV but its second life as ground truth, which has three
consequences.

First, evaluations and vendor claims of the form ``catches X\% of KEV''
measure agreement with CISA's visibility, which in our population covers at most
a third of observed exploitation. A method that finds exploited CVEs that KEV
has missed, the very use the white paper proposes for LEV, is penalized for
doing so.

Second, the strata translate into the cost of acting on LEV. Read as triage, a
defender working through CVEs with LEV of at least 0.9 examines about 4.6 CVEs
per exploited CVE found when KEV is the truth, but about 2.1 against the
estimated truth. In the middle stratum the cost is at most about 8 (from Chao's
bound), and in the lowest at most about 89. KEV-based figures more than double
the apparent cost of working the top of the list.

The same counts price the white paper's proposal to review high-LEV CVEs for
inclusion in KEV, a workload it expects to be
small~\cite[Sec.~3.3]{mell2025lev}. In our population, 498 CVEs have LEV of at
least 0.9 and are not in KEV. We estimate that about 168 of them were
exploited, one in three, and 139 of those are already recorded by SENSOR or
VULNCHK. A reviewer would therefore examine about three CVEs for each exploited
CVE that KEV lacks, and would find most of them sooner by reading the other
lists first. In the middle stratum, 1,895 CVEs are not in KEV and, by Chao's
bound, at least about 180 of them were exploited, so the cost is at most about
11 reviews for each. Whether a CVE found this way meets KEV's own criteria of
scope and remediation is a separate question.

Third, provenance determines what can be combined. Our dependence audit was
possible only because VulnCheck publishes the source of every report behind
each entry. Catalogues that publish per-entry provenance can be combined;
catalogues that do not can only be compared.

\section{Threats to validity}\label{sec:threats}

The repository records 55 threats with their mitigations. We group
those that bear most on the conclusions under the four usual headings.

\paragraph{Construct validity.} The lists measure different things. KEV lists
only CVEs with remediation guidance and federal relevance, and honeypots see
only exploitation of internet-reachable services. These structural zeros violate
the shared-population assumption of capture-recapture, stratifying by LEV does
not repair them, and they bias the estimates downwards; the over-confidence
figures are therefore upper bounds. Telegram sightings include claims we cannot
verify, so they are not a primary list; Table~\ref{tab:config} adds them back. 38\%
of VulnCheck-only CVEs rest on reports from one source with no sensor evidence;
the strict configuration, which requires a sensor or two independent sources,
preserves the top-stratum result.
Finally, \texttt{dateAdded} is a cataloguing date, so every timing comparison
with KEV is lagged by an unknown amount.

\paragraph{Internal validity.} Circularity runs in two directions: KEV is an
EPSS input, and Shadowserver is an EPSS label source. We measure the first
(Section~\ref{sec:res-leak}) and drop the second as a check
(Section~\ref{sec:res-cr}); neither can be removed from EPSS itself. Because
listing raises later EPSS, capture can move CVEs into the top stratum. Strata
formed on LEV that cannot see a CISA listing keep the top stratum's
over-confidence below about 2.2 against Chao's bound, but leave the log-linear
estimate unidentified (Section~\ref{sec:res-cr}).
The log-linear models assume no three-way interaction; we report estimates only
where the identification rules of Section~\ref{sec:m-cr} are met; those rules
were not registered in advance, and Table~\ref{tab:config} shows
how each configuration changes them. Two catalogues contain CISA KEV and
VulnCheck contains Shadowserver's list. Each was removed by source before
modelling, but a containment we did not detect would bias the estimates
downwards. Neither version normalization separates a corpus-wide shift from
genuine improvement in the newer model, so the artifact is reported as a range.
The backfill filter was added after a first reproduction miss
(Section~\ref{sec:res-repro}), although it is justified independently.

\paragraph{External validity.} SENSOR starts in October 2024, so the
capture-recapture population covers CVEs first scored over 15 months and
observed for only 6 to 21 months. Catchability changes with age: in CIRCL's
comparison of KEV-style catalogues, Shadowserver lists CVEs at a median age of
532 days against 252 for CISA~\cite{circl2026kevcatalogs}, so late listings are
still to come. Results for older CVEs, and for EPSS versions after v5, may
differ. Our search for prior work (arXiv, bibliographic indexes and the web, last
repeated on 15 November 2026) is structured but not a formal systematic review,
and may miss peer-reviewed work that those sources do not index.

\paragraph{Conclusion validity.} LEV's windows are 30 days, whereas our
snapshots are about 30.4 days apart. On 481 CVEs stratified by LEV decile, LEV
from true daily windows differs from ours by 3.9\% of mean LEV
(population-weighted), with no consistent sign. Analyses of RQ3 and RQ4 use
version-normalized LEV; as-specified LEV gives higher over-confidence
everywhere except the top stratum, where the two agree (Table~\ref{tab:ladder}
and Section~\ref{sec:res-cr}). VulnCheck's dates come from a later snapshot and
may be backdated, which affects its timing comparisons with CISA.

\section{Related work}\label{sec:related}

\paragraph{Exploit prediction.} Learning which vulnerabilities get exploited
predates EPSS. Allodi and Massacci compared severity scores with exploitation
in case-control studies~\cite{allodi2014casecontrol}; Sabottke et al.\ predicted
real-world exploits from Twitter~\cite{sabottke2015twitter}; Suciu et al.\
forecast the development of functional exploits~\cite{suciu2022expected}; and
Jacobs et al.\ used exploit prediction to inform remediation
policy~\cite{jacobs2020remediation}. EPSS was introduced by
Jacobs et al.~\cite{jacobs2021epss} and extended in version
3~\cite{jacobs2023epssv3}. Parla evaluates EPSS on high-severity CVEs in
KEV~\cite{parla2024}. Nucleus Security reports that EPSS
moves far more after KEV listing than before~\cite{netzer2026late}, consistent
with Section~\ref{sec:res-leak}. Mell et al.\ note that EPSS's training labels
come mainly from network intrusion-detection sensors, which can miss
exploitation that only other mechanisms see~\cite{mell2024weaknesses}.
Concurrently with our work, Lu and Stiawan audit EPSS's rank stability across
the v2--v5 releases and its effect on prioritization
workload~\cite{lu2026queueshock}. Their focus is the disruption to operational
queues at each release; ours is what the same discontinuities do to a metric
that multiplies scores across them. None of these works evaluates LEV.

\paragraph{LEV.} Beyond the white paper~\cite{mell2025lev}, an independent
reimplementation by the Risk Based Prioritization project questions LEV's
small-probability approximation and its independence
assumption~\cite{rbp2025lev}, but does not validate LEV against exploitation
data. The white paper itself lists a calibration error in EPSS, and the
amplification of dependent errors by the product, among the reasons why a
high-LEV CVE may be absent from KEV~\cite[Sec.~8]{mell2025lev}; it does not
measure either. LEV is already consumed downstream. Huang et al.\ use it, alongside EPSS
and KEV, as an exposure probability for incident mitigation in cyber-physical
systems~\cite{huang2026cps}, and newer prioritization work cites it as a
complement to EPSS without testing it~\cite{akhoundali2026gpm}. We found no
empirical evaluation of LEV.

\paragraph{KEV as ground truth.} Recent frameworks combine CVSS, EPSS and
KEV~\cite{shimizu2026chaining,podder2026cape}, build composite risk indicators
evaluated on KEV~\cite{sherif2026kri}, train a ranking model with KEV entries as
weak supervision~\cite{motamedi2026threatlens}, order dependency upgrades by
KEV membership~\cite{omondi2026kevgraph}, or validate an EPSS-based measure of
exploitation pressure by the share of KEV entries it captures~\cite{almarashi2026entropy}. Prospective evaluations increasingly
guard against temporal leakage in the evidence~\cite{alpay2026graphs}, but the
label itself is still taken as complete. The GitHub Popularity Metric is
evaluated against CISA KEV and against a pooled set that also includes
VulnCheck's KEV~\cite{akhoundali2026gpm}; our audit shows that the latter
contains the former by construction, so the pooled set is not an independent
second check. Shaffer and Voicu model exploitation hazard with a decay
calibrated on KEV timing~\cite{shaffer2026hazard}; our cohorts show EPSS
rising, not decaying. Industry analyses already show that KEV is incomplete and
lagging~\cite{roytman2026kevparadox,garrity2026state,beardsley2025kevology}, and
CIRCL has compared five KEV-style catalogues
descriptively~\cite{circl2026kevcatalogs}. We add a population estimate, the
dependence audit that such comparisons need, and the consequence of both for
one metric.

\paragraph{Learning from incomplete labels.} Exploitation lists provide
positives and unlabelled CVEs, the setting of positive-unlabelled
learning~\cite{elkan2008pu,bekker2020pu}. Under its
selected-completely-at-random assumption, ranking evaluations survive
incompleteness. Our dependence audit and the age effects in
Table~\ref{tab:age} indicate that the assumption fails for KEV.

\paragraph{Multiple-systems estimation.} Capture-recapture over administrative
lists is standard in epidemiology~\cite{iwgdmf1995} and in public
policy~\cite{bird2018mse}. It has been used to estimate casualties and modern
slavery, where its sensitivity to list dependence and to sparse,
non-overlapping cells is well documented~\cite{binette2022mse,chan2021sparse}.
In software engineering it estimates the defects that inspections
miss~\cite{petersson2004inspection}. Closest to our setting, Meurs et al.\ link
police reports, incident-response cases and leak sites to estimate the number
of ransomware attacks on Dutch businesses~\cite{meurs2026ransomware}. Overlap
between independent discoveries has also been used to estimate vulnerability
rediscovery rates~\cite{herr2017rediscovery,ablon2017zeroday}. We found no
prior application to vulnerability exploitation.

\section{Conclusion}\label{sec:conclusion}

Evaluated against KEV, LEV looks badly over-confident, but much of that is not
LEV's doing. EPSS model releases put an artifact into the product that LEV
computes; EPSS scores for unlisted CVEs do not decay, so LEV accumulates without
the saturation its design assumes; and KEV records a minority of the
exploitation that other evidence-citing sources observe. Once these effects are
accounted for, LEV's most confident scores are over-confident against
estimated true exploitation by about a factor of two, rather than the 4.5 that
KEV suggests.
Below that, the data bound the error without identifying it: at most 2.5 in the
middle stratum, and possibly none in the lowest. A monotone correction brings
LEV's values in line with the broader ground truth a year ahead, though it must
be refitted as CVEs age, and LEV cannot anticipate most listings.

The lowest LEV stratum holds most CVEs, and four lists cannot identify it. A
fifth list with independent sensor coverage, such as GreyNoise's, is the
natural next step. The method is not specific to LEV: any metric validated
against KEV can be re-scored against the same truth ladder, using the per-CVE
truth bounds we release.

\texttt{make fetch} downloads the inputs that are not committed with the code
(EPSS, the CVE List and, with a token, VulnCheck's current feed),
\texttt{make verify-data} re-checks the cached EPSS files against their
recorded hashes, and \texttt{make reproduce} regenerates every table and
figure offline. Results that use VulnCheck reproduce exactly only from our
VulnCheck snapshot, whose hash is recorded; a newer feed raises a warning and
gives different numbers. SHA-256 digests, release tags and
retrieval dates for all inputs are listed in the repository. The truth bounds
give one row per CVE in the capture-recapture population: LEV as specified and
version-normalized, membership of each list, and the probability of exploitation
under counting, under Chao's bound, and under the log-linear estimate where it
is identified. A variant built from CISA, SENSOR and SOCIAL alone contains no
VulnCheck-derived data and will be released with the code; the primary-list version will be released if VulnCheck's
terms permit.

\section*{Ethics}

All inputs are public or available under a free community licence, and no
personal data are processed. Telegram-derived sightings are used only as
aggregated counts per CVE, as CIRCL publishes them. The truth bounds hold CVE
identifiers, LEV values, list membership and probabilities, and name no person
or organization.

\section*{Data availability}\label{sec:availability}

The code, the pinned inputs whose licences allow redistribution, and the
per-CVE truth bounds are available at
\url{https://github.com/kumarswamyg2005/lev-eval}. \texttt{make fetch}
downloads the inputs that are not committed with the code (EPSS, the CVE List
and, with a free community token, VulnCheck's current feed), and \texttt{make
reproduce} regenerates every table and figure. VulnCheck's snapshot cannot be
shared, because its terms do not state that the data may be redistributed.
Results that use it, including all primary-list results, therefore reproduce
exactly only from our snapshot, whose hash is recorded; VulnCheck serves only
its current feed, so a re-run gives slightly different numbers. For the same
reason the released truth bounds are the variant built from CISA, SENSOR and
SOCIAL alone; the primary-list version will be released if VulnCheck's terms
permit.

\section*{CRediT authorship contribution statement}

\textbf{Gurram Naga Kumaraswamy:} Conceptualization, Methodology, Software,
Formal analysis, Data curation, Visualization, Writing -- original draft.
\textbf{Kakarla Saikrishna:} Conceptualization, Methodology, Validation,
Investigation, Writing -- original draft.
\textbf{Bulla Rajesh:} Supervision, Writing -- review and editing.

\section*{Funding}

This research did not receive any specific grant from funding agencies in the
public, commercial, or not-for-profit sectors.

\section*{Declaration of competing interest}

The authors declare that they have no competing interests.

\section*{Declaration of generative AI and AI-assisted technologies in the manuscript preparation process}

During the preparation of this work the authors used Claude (Anthropic) to
assist with drafting and editing the text, with preparing the \LaTeX{} source,
and with writing and reviewing the analysis code. The research questions, the
study design, the choice of data and methods, and the interpretation of the
results are the authors' own. After using this tool, the authors checked every
reported number against the output of the released code, verified each
reference against its source, reviewed and edited the content as needed, and
take full responsibility for the content of the published article.
